\documentclass[12pt]{article}
\usepackage[utf8]{inputenc}
\usepackage[T1]{fontenc}
\usepackage[margin=1in]{geometry}
\usepackage{setspace}
\usepackage{xcolor}
\usepackage[colorlinks=true, urlcolor=blue!60!black, linkcolor=blue!60!black]{hyperref}
\setlength{\parindent}{2em}
\setlength{\parskip}{1.25em}
\renewcommand{\baselinestretch}{1.0}
\begin{document}
\begin{center}
{\LARGE \textit{Chaplinsky v. New Hampshire}, 315 U.S. 568 (1942)\par}
\vspace{0.8em}
{\large Case Brief --- Four-Page Edition\par}
\end{center}
\vspace{1em}

\section{Memory Jogger}

A Jehovah's Witness called a city marshal a ``damned racketeer'' and ``damned Fascist'' on a public street; a unanimous Court upheld his conviction, holding that ``fighting words,'' face-to-face epithets likely to provoke an average person to immediate violence, belong to narrow classes of speech the First Amendment does not protect (\href{https://supreme.justia.com/cases/federal/us/315/568/}{Justia}).

\section{Detailed Case Facts}

On a busy Saturday afternoon in Rochester, New Hampshire, Walter Chaplinsky, a Jehovah's Witness, distributed religious literature on a public sidewalk and denounced organized religion as a ``racket.'' Citizens complained to City Marshal Bowering, who told them Chaplinsky was acting lawfully but warned Chaplinsky that the crowd was growing restless. After a disturbance broke out, a traffic officer escorted Chaplinsky toward the police station without arresting him.

On the way, near City Hall, they met Bowering hurrying toward the reported riot. After Bowering repeated his warning (and, Chaplinsky testified, cursed him when he demanded the arrest of the crowd), Chaplinsky called Bowering ``a God damned racketeer'' and ``a damned Fascist,'' adding that ``the whole government of Rochester are Fascists or agents of Fascists.'' Chaplinsky admitted the words except for invoking the name of God.

\section{Procedural History}

Chaplinsky was convicted in municipal court under a New Hampshire law providing that ``[n]o person shall address any offensive, derisive or annoying word to any other person who is lawfully in any street or other public place.'' On appeal, a jury convicted him after a trial \textit{de novo}; the court excluded evidence of his religious mission, his mistreatment by the crowd, and police neglect. The New Hampshire Supreme Court affirmed, construing the statute to reach only face-to-face ``fighting words'' likely to provoke an average addressee to violence. The U.S. Supreme Court affirmed on March 9, 1942.

\section{Judicial Votes}

Affirmed, \textbf{9--0}.
\begin{itemize}
\item \textbf{Majority:} Justice Murphy, for a unanimous Court (Chief Justice Stone and Justices Roberts, Black, Reed, Frankfurter, Douglas, Byrnes, and Jackson).
\item \textbf{Concurrences and dissents:} None.
\end{itemize}

\section{Holding}

Certain ``well-defined and narrowly limited classes of speech,'' including ``fighting words'' that ``by their very utterance inflict injury or tend to incite an immediate breach of the peace,'' have such slight social value as a step to truth that any benefit is outweighed by the social interest in order and morality. As narrowly construed and applied to Chaplinsky's epithets, the statute is neither unconstitutionally vague nor an impermissible abridgment of free speech.

\section{Analysis of Opinions}

\textbf{Majority (Murphy): scope.} Only free speech was really at stake: the case involved spoken words, not the press, and ``cursing a public officer'' is not an exercise of religion.

\textbf{Unprotected categories.} Free speech is ``not absolute at all times and under all circumstances.'' Drawing on Zechariah Chafee, Murphy listed classes whose punishment had ``never been thought to raise any Constitutional problem'': the lewd and obscene, the profane, the libelous, and fighting words. Such utterances are ``no essential part of any exposition of ideas.''

\textbf{The state court's construction.} The Court accepted the New Hampshire court's narrowing of the statute to face-to-face words having ``a direct tendency to cause acts of violence by the persons to whom, individually, the remark is addressed,'' judged objectively by what ``men of common intelligence'' would understand.

\textbf{Vagueness and application.} So narrowed, the statute punished specific ``verbal acts'' and gave fair notice. Taking judicial notice that the epithets were likely to provoke retaliation, the Court upheld the conviction and held that provocation and truth need not be admitted as defenses.

\textbf{Later developments.} \textit{Cohen v. California} (1971) limited fighting words to ``personally abusive epithets'' directed face to face; \textit{Gooding v. Wilson} (1972) required an \textit{immediate} breach of the peace, effectively discarding the ``inflict injury'' prong; Powell's \textit{Lewis} (1974) concurrence suggested police must show more restraint. \textit{R.A.V.} (1992) barred viewpoint discrimination within the category, and \textit{United States v. Stevens} (2010) rejected ad hoc value-balancing as ``startling and dangerous,'' reading \textit{Chaplinsky} as describing ``historic and traditional categories long familiar to the bar.''

\section{Examples: Future Applications}

\begin{enumerate}
\item \textbf{Same-Side (Upheld): Face-to-Face Challenge.} At a protest, an activist inches from a private counter-protester screams personal profanities and a challenge to fight, under a narrowly construed, content-neutral statute.
\item \textbf{Same-Side (Upheld): Slurs at a Neighbor's Door.} Personal slurs inches from a neighbor's face, with a dare to fight, are personally abusive epithets likely to provoke violence.
\item \textbf{Opposite-Side (Protected): Bullhorn and Online.} ``Fascist pigs!'' shouted across the street at police, or posted online, lacks a face-to-face risk of immediate retaliation.
\item \textbf{Opposite-Side (Invalid): Bias-Only Ban.} An ordinance punishing only race-, religion-, or gender-based fighting words is viewpoint discrimination under \textit{R.A.V.}
\item \textbf{Fence-Sitter: Chaplinsky's Words Today.} Calling an arresting officer a damned racketeer and a damned Fascist, without threat, is face to face yet political criticism of an official, and the Court has not upheld a fighting-words conviction since.
\end{enumerate}

\section{Critique}

\textbf{Ironic facts.} Chaplinsky was criticizing an official over police failure to protect an unpopular religious minority, core political speech that \textit{New York Times v. Sullivan}, \textit{Cohen}, and \textit{City of Houston v. Hill} now protect.

\textbf{Outdated assumptions.} The assumption that ``damned Fascist'' automatically triggers violence reflects a gender- and class-bound ``average addressee'' honor code inviting selective enforcement. And the ``inflict injury'' prong collides with \textit{Snyder v. Phelps} and \textit{Matal v. Tam}, which hold that emotional distress alone cannot strip speech of protection.

\section{Key Quotations}

\begin{quote}
``These include the lewd and obscene, the profane, the libelous, and the insulting or `fighting' words---those which by their very utterance inflict injury or tend to incite an immediate breach of the peace. It has been well observed that such utterances are no essential part of any exposition of ideas, and are of such slight social value as a step to truth that any benefit that may be derived from them is clearly outweighed by the social interest in order and morality.''
\end{quote}
\noindent --- Justice Murphy, for the Court, 315 U.S. at 572.

\begin{quote}
``Allowing the broadest scope to the language and purpose of the Fourteenth Amendment, it is well understood that the right of free speech is not absolute at all times and under all circumstances. There are certain well-defined and narrowly limited classes of speech, the prevention and punishment of which have never been thought to raise any Constitutional problem.''
\end{quote}
\noindent --- Justice Murphy, for the Court, 315 U.S. at 571--572.

\begin{quote}
``The test is what men of common intelligence would understand would be words likely to cause an average addressee to fight.''
\end{quote}
\noindent --- Justice Murphy (quoting the New Hampshire Supreme Court), 315 U.S. at 573.

\end{document}
